% Consolidated textbook-ch09 from f43cd09215a13d29f88baa670fa48441ec2a44e8; source fragments preserved.
% 2026-11-24 八批严格复核第3批：旧片段哈希对应历史版，当前文件含二校修订。

% === BEGIN textbook-ch09-choices.tex blob c8a64a6b89aa4da59d76e76b1707e1bc3b48ae72 ===
% 用户扫描《大学物理简明教程》第5版第9章“变化的电磁场”
% 选择题9.1.1—9.1.13，印刷/PDF第236—237页。
% 题意研究性转述，推导独立；第9.1.11四选项按严格定义均有疑点。

\subsection{选择题：电磁感应、互感及位移电流}

\sourceproblem{9.1.1}{236}{236}
\problemtext 一圆形金属线圈置于均匀且恒定的磁场中，
问下列运动哪些产生感应电流：平移、绕与磁场平行或垂直的直径转动等。
\workings
闭合回路中的感应电动势满足
\[
\mathcal E=-\frac{d}{dt}(\bm B\cdot\bm S)
=-BS\frac d{dt}\cos\theta
\]
（$B,S$恒定）。
纯平移、绕与$\bm B$平行的直径转动都不改变磁通量。
若绕与$\bm B$垂直的直径转动，
面法线与磁场夹角随时间变化，一般有$\mathcal E\ne0$。
故选\textbf{B}。
\crosscheck 不应因“闭合导体有运动”就断定有感应电流：
关键是闭合环路的总磁通量是否随时间改变。

\sourceproblem{9.1.2}{236}{236}
\problemtext 图9.1.2中一根长度$l$的直导线在均匀磁场$\bm B$中以
速度$\bm v$移动，$\bm v$、导线与$\bm B$位于同一纸面，
分别与水平成一定夹角。问导线端点间的动生电动势。
\workings
对刚性平移导体段，
\[
\mathcal E_{ab}=\int_a^b(\bm v\times\bm B)\cdot d\bm l.
\]
$\bm v$和$\bm B$都位于纸面，
故$\bm v\times\bm B$垂直纸面，而导线元$d\bm l$在纸面内，
其点积处处为零。因此
\[
\boxed{\mathcal E_{ab}=0},
\]
原书选\textbf{D}。
\crosscheck 不能把单个夹角$\alpha$生搬进$Blv\sin\alpha$；
该特殊公式还要求明确磁场、导线和速度的三维取向。

\sourceproblem{9.1.3}{236}{236}
\problemtext 关于涡旋电场的四项描述中，
判断“不正确”的一项：涡旋场能对电荷施力、来自变化的磁场、
由静电荷激发、其场线可闭合。
\workings
法拉第--麦克斯韦方程
\[
\nabla\times\bm E_{\mathrm{ind}}
=-\frac{\partial\bm B}{\partial t}
\]
说明磁场的时间变化激发具有非零环流的感应电场，
它对电荷仍有力$q\bm E_{\mathrm{ind}}$，
感应场线在典型对称情形下形成闭合环。
静止电荷激发的库仑静电场是保守场，并非此处所说的涡旋感应场。
故错误项是\textbf{C}。
\crosscheck 对于动态的一般电场可分解静电样部分与非保守部分；
不能把“涡旋电场”误认为“变化电场”本身。

\sourceproblem{9.1.4}{236}{236}
\problemtext 静电场、稳恒磁场、感生电场、随时间变化磁场四类矢量场中，
哪一类是势函数描述下的保守场？
\workings
静电场满足$\nabla\times\bm E_{\mathrm{static}}=\bm0$，
在无奇异、单连通区域可表示为
$\bm E_{\mathrm{static}}=-\nabla\varphi$，
使任意闭路$\oint\bm E_{\mathrm{static}}\cdot d\bm l=0$。
感生电场一般环流非零；磁场本身不按保守力场定义。
故选\textbf{A}。
\crosscheck 此处“保守场”指电场作为单位电荷受力场，
而非仅凭磁场是否随时间改变来判别。

\sourceproblem{9.1.5}{236}{236}
\problemtext 一圆形铜线环和形状完全相同但材料不同的不锈钢线环静置，
穿过两者的磁通量变化率相同，忽略自感，比较感应电动势大小。
\workings
两环均满足法拉第定律
$\mathcal E=-d\Phi_B/dt$，与材料电阻率无关，
因此
\[
\boxed{|\mathcal E_{\mathrm{Cu}}|
=|\mathcal E_{\mathrm{steel}}|},
\]
原书选\textbf{D}。
\crosscheck 材料电阻不同会使感应电流$I=\mathcal E/R$不同，
但题问的是电动势而不是电流。

\sourceproblem{9.1.6}{236}{236}
\problemtext 一匀强磁场中转轴与$\bm B$垂直的线圈匀速转动，
转速增为原来的2倍、其他条件不变。比较最大感应电动势、
电流和交流频率。
\workings
在单匝或固定匝数$N$的转动线圈中，
$\Phi=NBA\cos(\omega t+\phi)$，
\[
\mathcal E=NBA\omega\sin(\omega t+\phi).
\]
最大感应电动势$\mathcal E_{\max}=NBA\omega$，
故转速加倍时最大值加倍，原题选\textbf{A}。
\crosscheck 频率$f=\omega/(2\pi)$同样仅加倍而非四倍；
电流峰值还受电路电阻、电感和频率依赖的阻抗影响；纯电阻且阻值固定时电流峰值也加倍，一般阻抗下不能无条件断言具体倍数。

\sourceproblem{9.1.7}{236}{236}
\problemtext 长$L$的金属杆绕端点$O$以角速度$\omega$在垂直恒定磁场
$\bm B$的平面内匀速转动，求两端动生电动势。
\workings
距轴$r$的杆微元切向速率$v=\omega r$，
局部磁力单位电荷沿杆方向，故
\[
|\mathcal E|=\int_0^L B\omega r\,dr
=\boxed{\frac12B\omega L^2}.
\]
原题选\textbf{E}。
\crosscheck 动生电动势与时间无关；
杆角度$\omega t+\theta$只决定杆在平面中的瞬时朝向，
并不使匀强垂直磁场下的电动势额外乘余弦。

\sourceproblem{9.1.8}{237}{237}
\problemtext 两回路互感分别记为$M_{21}$与$M_{12}$，
且$|di_1/dt|<|di_2/dt|$。比较互感系数及线圈2中由线圈1产生的
电动势$\varepsilon_{21}$与线圈1中由线圈2产生的电动势$\varepsilon_{12}$的大小。
\workings
对线性静止互易电路，互感系数满足
\[
\boxed{M_{12}=M_{21}=M}.
\]
而互感电动势的大小分别为
$|\varepsilon_{21}|=M|di_1/dt|$、
$|\varepsilon_{12}|=M|di_2/dt|$，
若$M\ne0$，故
\[
\boxed{|\varepsilon_{21}|<|\varepsilon_{12}|},
\]
原书选\textbf{D}。
\crosscheck 若$M=0$，两互感电动势均为零；一般公式中的幅值系数应取$|M|$。电动势比较的是绝对值；
有向瞬时电动势符号还取决于线圈绕向、端点参考方向与电流变化符号。

\sourceproblem{9.1.9}{237}{237}
\problemtext 磁能公式$W_m=\tfrac12LI^2$的适用范围，
是否仅限无限长螺线管、单匝线圈、密绕螺线环，或适用于自感恒定的任意线性回路？
\workings
线性磁场能量积累时
$dW_m=I\,d(LI)=LI\,dI$，
在$L$不随电流变化时积分给
\[
\boxed{W_m=\int_0^I L i\,di=\frac12LI^2}.
\]
故适用于自感系数$L$固定的任何线性几何回路，
原题选\textbf{D}。
\crosscheck 对铁磁非线性材料，$L$随$I$可能变化，
一般需积分$W=\int I\,d\Psi$，不能直接写$\tfrac12LI^2$。

\sourceproblem{9.1.10}{237}{237}
\problemtext 真空中一根无限长直导线载流$I$，
问距导线$a$处的磁场能量密度。
\workings
导线场强
$B(a)=\mu_0I/(2\pi a)$，
真空磁能密度$w_m=B^2/(2\mu_0)$，
因此
\[
\boxed{w_m=\frac1{2\mu_0}
\left(\frac{\mu_0I}{2\pi a}\right)^2
=\frac{\mu_0I^2}{8\pi^2a^2}}.
\]
原书选\textbf{B}。
\crosscheck 量纲$[\mathrm{J/m^3}]$；
不要把场能密度错写为$B^2/(2\varepsilon_0)$。

\sourceproblem{9.1.11}{237}{237}
\problemtext 判断位移电流的定义、起因、热效应和磁效应的四种断言：
(A)位移电流就是变化的电场；(B)由变化的磁场产生；
(C)热效应服从焦耳定律；(D)磁效应不服从安培环路定理。
\workings
定义为位移电流密度
\[
\boxed{\bm J_D=\frac{\partial\bm D}{\partial t}}
\]
（真空为$\varepsilon_0\partial\bm E/\partial t$），
其\emph{量纲和物理意义}不是“变化的电场”本身；
时间变化的电场还需乘以$\varepsilon_0$并对时间求导。
变化的磁场按法拉第方程激发涡旋电场，
但不是位移电流\emph{定义}中的必要起因。
理想真空位移电流没有自由载流子欧姆耗散，
不遵循“$I^2R$”的焦耳定律；
它却要计入推广后的安培--麦克斯韦环路定理
\[
\oint_L\bm H\cdot d\bm l
=I_{\mathrm{cond}}+\frac{d}{dt}\int_S\bm D\cdot d\bm S.
\]
所以按四个选项的\emph{严格字面定义}，
\textbf{均不能无条件成立，原题缺少准确选项}。
\crosscheck \textbf{原题待二校：}若教材将(A)宽泛简称为
“位移电流来自变化的电场”，可理解其想选(A)，
但其字面没有量纲一致性，不应把不严谨断言冒充严格定义。

\sourceproblem{9.1.12}{237}{237}
\problemtext 一个充好电的平行板电容器，
左极板$A$正、右极板$B$负，外电路电阻及开关如图9.1.12。
闭合开关后，板间位移电流沿图示$x$轴的哪个方向？
\workings
闭合开关后电容器放电，正负极板电量绝对值逐渐减小。
电场$\bm E$初始从$A$指向$B$，即$+x$方向，
其大小递减，故
$\partial\bm D/\partial t$沿$-x$方向。
因此板间位移电流密度满足
\[
\boxed{\bm J_D=\frac{\partial\bm D}{\partial t},\qquad J_{D,x}<0}.
\]
方向为\textbf{$x$轴负向}，对应原题选\textbf{B}。
\crosscheck 充电时位移电流与电场增长同向；
本题是\emph{放电}，容易误将充电时的方向套用到放电。

\sourceproblem{9.1.13}{237}{237}
\problemtext 平行板电容器\emph{充电}时，
画两条闭合环路$L_1,L_2$，
分别穿过板间的位移电流和引线中的传导电流，
比较$\oint_{L_1}\bm H\cdot d\bm l$、
$\oint_{L_2}\bm H\cdot d\bm l$。
\workings
使用安培--麦克斯韦定律
\[
\oint_L\bm H\cdot d\bm l
=I_{\mathrm{cond,encl}}+I_{\mathrm{disp,encl}}.
\]
同一传输回路充电时，两种电流的\emph{有向数值}大小相等，
$I_D=\varepsilon_0\,d\Phi_E/dt=I$，
只要$L_1,L_2$按同一右手定向约定，
\[
\boxed{\oint_{L_1}\bm H\cdot d\bm l
=\oint_{L_2}\bm H\cdot d\bm l=I}.
\]
原题选\textbf{B}。
\crosscheck 不能只对一条曲面写传导电流而把另一条曲面的位移电流漏掉；
这正是推广安培环路定理的原因。
% === END textbook-ch09-choices.tex ===

% === BEGIN textbook-ch09-blanks.tex blob 3bf3c0ff0a9ae712a2a053354eef68f63979c9b9 ===
% 《大学物理简明教程》第5版第9章填空题9.2.1—9.2.14
% 原书/PDF第238—239页；据用户扫描页校订题意并独立推导。

\subsection{填空题：感应电动势、位移电流及磁场能量}

\sourceproblem{9.2.1}{238}{238}
\problemtext 把一金属圆环从磁极间沿垂直于磁感应强度的方向抽出，
问圆环会受到哪种电磁作用。
\workings
闭合环路在移出磁场边界时，穿过回路的磁通量随时间变化，
$\mathcal E=-d\Phi_B/dt$，由此产生感应电流。
感应电流在磁场内受安培力，按楞次定律该合力阻碍移出运动。
因此填空为\textbf{磁阻力（阻碍抽出的安培力）}。
\crosscheck 只有圆环\emph{部分跨越磁场边界}才有非零磁通量变化；
若整环在空间均匀、无限延伸的磁场中匀速平移，则总电动势为零。

\sourceproblem{9.2.2}{238}{238}
\problemtext 分别填出动生电动势、感生电动势对应的非静电场力来源，
以及产生感生电场的场源。
\workings
动生电动势来自运动电荷所受磁力
$\boxed{\bm F_{\mathrm{mot}}=q\bm v\times\bm B}$；
感生电动势来自涡旋电场对电荷的力
$\boxed{\bm F_{\mathrm{ind}}=q\bm E_{\mathrm{ind}}}$；
感生电场的场源是\textbf{随时间变化的磁场}，满足
\[
\boxed{\nabla\times\bm E_{\mathrm{ind}}
=-\frac{\partial\bm B}{\partial t}}.
\]
\crosscheck 电磁感应有动生和感生两类物理机制，
不可全部写作洛伦兹磁力或全部写作静电库仑力。

\sourceproblem{9.2.3}{238}{238}
\problemtext 长$\ell$金属直导线在垂直匀强磁场$\bm B$的平面内
绕固定轴以角速度$\omega$转动。
转轴取在导线的哪里时两端电动势最大，最大值是多少；
取在何处两端电动势最小，最小值是多少？
\workings
设转轴距导线一端为$a$，
从轴向该端和另一端分别有距离$a,\ell-a$，
有向端电势差
\[
\mathcal E=\frac12B\omega
\bigl[(\ell-a)^2-a^2\bigr]
=\frac12B\omega(\ell^2-2\ell a).
\]
若转轴限制在杆上，即$0\le a\le\ell$，在$a=0$或$a=\ell$时，
\[
\boxed{|\mathcal E|_{\max}=\frac12B\omega\ell^2},
\]
转轴位于\textbf{导线端点}；
当$a=\ell/2$即转轴位于\textbf{中点}时，
\[
\boxed{|\mathcal E|_{\min}=0}.
\]
\crosscheck 两半段虽各自具有对轴的动生电动势，
但关于轴的有向贡献相反，故绕中点转动时\emph{端到端}电动势零。若允许转轴在杆延长线上，$|\mathcal E|$可随轴距增大，理想模型中无有限全局最大值。

\sourceproblem{9.2.4}{238}{238}
\problemtext 通过单匝线圈的磁通量为
$\Phi(t)=6t^2+9t+2$（SI单位），求$t=2$ s的感应电动势。
\workings
法拉第定律给
\[
\mathcal E(t)=-\frac{d\Phi}{dt}=-(12t+9)\,\mathrm V,
\quad\boxed{\mathcal E(2\,\mathrm s)=-33\,\mathrm V}.
\]
若只求电动势\emph{大小}则是$\boxed{33\,\mathrm V}$。
\crosscheck 负号是相对于预设线圈正绕行方向的定向符号；
原题没有给明确的绕行方向时应同时说明大小和有向取值约定。

\sourceproblem{9.2.5}{238}{238}
\problemtext 写出法拉第电磁感应定律的闭合积分表达式。
\workings
静止闭合路径的场方程为
\[
\boxed{\oint_{\partial S}\bm E\cdot d\bm l
=-\frac d{dt}\int_S\bm B\cdot d\bm S
=-\frac{d\Phi_B}{dt}}.
\]
\crosscheck 环路和面的正向应满足右手关系；
若边界随时间运动，还必须将运动边界的
$\bm v\times\bm B$项计入总电动势。

\sourceproblem{9.2.6}{238}{238}
\problemtext 写出时间变化磁场与它所激发的涡旋电场的微分关系。
\workings
麦克斯韦--法拉第方程为
\[
\boxed{\nabla\times\bm E=-\frac{\partial\bm B}{\partial t}}.
\]
\crosscheck 左边单位为$\mathrm{V/m^2}$，
右边单位为$\mathrm{T/s}$，两者在SI制下等价；
这也说明一般感生电场不满足无旋静电场条件。

\sourceproblem{9.2.7}{238}{238}
\problemtext 无限长直导线竖直放置、电流向上；
在其右侧有同平面的任意形状闭合线圈。
(1) 线圈整体沿导线向下运动；
(2) 线圈在本平面内以垂直导线的速度向导线靠近。
约定\emph{顺时针}感应电动势为正，分别判定正负。
\workings
(1) 无限长稳恒电流导线的磁场在竖直方向平移不变，
线圈平行移动时磁通量不变：
\[
\boxed{\mathcal E_1=0}.
\]
(2) 右侧磁场指向纸内，
线圈靠近导线使进入纸面的通量\emph{增大}，
感应电流应产生指向纸外的反向场，故沿\emph{逆时针}。
题设顺时针为正，所以
\[
\boxed{\mathcal E_2<0}.
\]
\crosscheck 特别区分约定：若规定逆时针为正，
第二项的\emph{数值符号}会反转，
但实际绕行方向不会改变。

\sourceproblem{9.2.8}{238}{238}
\problemtext 一$N$匝矩形线圈固定于小振动杆，
其宽度（垂直运动方向）为$b$，有一部分位于恒定垂直磁场$B$内。
杆的位置为$x(t)=A\cos\omega t$，
求忽略边缘效应时线圈中的感应电动势。
\workings
以边缘内的重叠长度$x$表示有效面积，
可取总磁链$\Psi_B=NBb\,x$（若图中$x$为向外位移则相应变号）。
因此
\[
\boxed{\mathcal E(t)=-NBb\dot x
=NBb\,A\omega\sin\omega t}.
\]
最大大小为$\boxed{\mathcal E_{\max}=NBbA\omega}$。
\crosscheck 这一式要求线圈的有源边一直跨越均匀磁场的\emph{边界}，
且另一边不同时跨界；否则有效面积函数不是简单的$b\,x$。

\sourceproblem{9.2.9}{238}{238}
\problemtext 在两条平行金属导轨上有可自由滑动的导体杆
$AB$（左）及$CD$（右），穿过纸面的均匀磁场向里，
外力使左杆$AB$向左运动。
问右杆$CD$将如何运动、磁场对左杆施加的合力指向何处。
\workings
$AB$向左移动使两杆组成的回路面积\emph{增加}，
向纸内的磁通量增加。
楞次定律要求感应电流产生向纸外的磁场，
因此整个回路按逆时针方向通电：
左杆中电流由$A$流向$B$（向下），
右杆$CD$由$D$流向$C$（向上）。
由$\bm F=I\bm l\times\bm B$，
右杆受力向\textbf{左}，左杆受力向\textbf{右}。
\crosscheck 左杆受磁力总与外力驱动的左移方向相反，
符合感应电流消耗机械输入能量的楞次定律。

\sourceproblem{9.2.10}{238}{238}
\problemtext 半径$r=20$ cm的圆形闭合回路位于垂直磁场
$B=1.0$ T中，半径从$t=0$以恒定速率
$dr/dt=-50$ cm/s缩小。
求初始感应电动势大小，以及让电动势保持该值时
需要的恒定面积变化率。
\workings
磁通量$\Phi=B\pi r^2$，
故
\[
|\mathcal E|=B\left|2\pi r\frac{dr}{dt}\right|
=1\cdot2\pi(0.20)(0.50)
=\boxed{0.20\pi\,\mathrm V\approx0.628\,\mathrm V}.
\]
若$B$不变，则要求
\[
\boxed{\frac{dS}{dt}=-\frac{|\mathcal E|}{B}
=-0.20\pi\,\mathrm{m^2/s}}.
\]
\crosscheck 保持电动势恒定意味着$S(t)$线性减小；
半径的收缩速率需随$r$减小而增大，
不能一直保持$dr/dt=-0.5$ m/s。恒定面积变化率只能维持到$S(t)=0$，在题设初始半径下为$t=0.20$ s，零面积处该收缩模型终止。

\sourceproblem{9.2.11}{238}{238}
\problemtext 细长螺线管横截面积$S=1\,\mathrm{cm^2}$，
总匝数$N=400$，电流$I=4$ A时内部均匀场$B=0.2$ T。
忽略漏磁，求自感。
\workings
总磁链$\Psi=NB S$，
在线性磁路、自感不随电流改变且无外加磁链时，$L=\Psi/I$，因此
\[
\boxed{L=\frac{NBS}{I}
=\frac{400\cdot0.2\cdot10^{-4}}4
=2.0\times10^{-3}\,\mathrm H}.
\]
\crosscheck $1\,\mathrm{cm^2}=10^{-4}\,\mathrm{m^2}$；
此处$B$已给，不必另外知道螺线管长度。

\sourceproblem{9.2.12}{238}{238}
\problemtext 任意两个固定几何形状的线性互感回路1、2，
流过相同的稳恒电流$I$。
由回路2电流穿过回路1的磁通量记$\Phi_{12}$，
反之记$\Phi_{21}$。求二者关系。
\workings
互感互易性给
$M_{12}=M_{21}=M$，
在线性静止互易介质及一致绕向约定下，若每个回路都是单匝或$\Phi$已定义为相应磁链，
则
\[
\Phi_{12}=M_{12}I,\qquad
\Phi_{21}=M_{21}I,
\]
故
\[
\boxed{\Phi_{12}=\Phi_{21}}.
\]
\crosscheck 多匝回路需明确比较的是总\emph{磁链}
而非随意某一匝的单匝磁通量。

\sourceproblem{9.2.13}{238}{238}
\problemtext 某一线圈电流在$0.002$ s内均匀地由$10$ A增至$12$ A，
其间自感电动势大小为400 V。
求该线圈的自感系数$L$。
\workings
线圈自感电动势
$|\mathcal E_L|=L|\Delta I/\Delta t|$，
故
\[
\boxed{L=\frac{400}{(12-10)/0.002}
=\frac{400}{1000}=0.40\,\mathrm H}.
\]
\crosscheck $[\mathrm V/(\mathrm{A/s})]=\mathrm H$；
电流增大时有向自感电动势取阻碍电流增长的方向。

\sourceproblem{9.2.14}{239}{239}
\problemtext 在真空中长直螺线管通电$I_1$时储磁能$W_1$。
若插入相对磁导率$\mu_r=4$的线性磁介质，
电流增为$I_2=2I_1$且几何不变，求$W_1:W_2$。
\workings
固定线圈匝数与几何，
线性自感$L\propto\mu$，
磁能$W=\tfrac12LI^2$，故
\[
\frac{W_2}{W_1}
=\frac{L_2}{L_1}\left(\frac{I_2}{I_1}\right)^2
=4\cdot2^2=16.
\]
\[
\boxed{W_1:W_2=1:16}.
\]
\crosscheck 此处假定磁介质的$\mu_r=4$为\emph{线性常数}；
对有磁滞的铁磁体并不一般满足能量恰按$\mu I^2$变化。
% === END textbook-ch09-blanks.tex ===

% === BEGIN textbook-ch09-solutions-001-006.tex blob 7492b9039bf66092f023435d113c5a15f5b754c4 ===
% 《大学物理简明教程》5版，第9章解答题9.3.1—9.3.6
% 用户扫描原书/PDF第239页，高分辨率核对几何、方向和数值；
% 独立研究性解析，所有题仍待专家二校。

\subsection{解答题：感生、动生电动势与穿越磁场边界}

\sourceproblem{9.3.1}{239}{239}
\problemtext 匀强磁场$B=0.4\,\mathrm T$，
圆形闭合回路垂直磁场，面积随时间变化为
$S(t)=5t^2+3$（单位$\mathrm{cm^2}$、$t$单位秒）。
求$t=2$ s时的感应电动势大小。
\workings
穿过回路的磁通量为
$\Phi=BS=0.4(5t^2+3)\times10^{-4}\,\mathrm{Wb}$。
法拉第定律给
\[
|\mathcal E|=\left|\frac{d\Phi}{dt}\right|
=B\,\frac{dS}{dt}
=0.4(10t)\times10^{-4}\,\mathrm V.
\]
$t=2$ s时
\[
\boxed{|\mathcal E|=8.0\times10^{-4}\,\mathrm V}.
\]
\crosscheck $1\,\mathrm{cm^2}=10^{-4}\,\mathrm{m^2}$；
若遗漏面积单位换算将导致$10^4$倍错误。

\sourceproblem{9.3.2}{239}{239}
\problemtext 长直导线载流$I$、电流按图向下，
其右侧同平面有半圆形金属杆$MeN$，半径$b$，
圆心$O$距直导线$a$且$a>b$。
半圆杆整体以速率$v$沿直导线向上匀速移动。
求杆中动生电动势的大小和端点电势差$U_M-U_N$。
\workings
设从长直导线向右量$x$，
杆端横坐标$x_M=a-b,x_N=a+b$。
由于直导线向下通电，
杆位置处磁场垂直纸面\emph{向外}，大小
\[
B(x)=\frac{\mu_0I}{2\pi x}.
\]
速度向上，故$\bm v\times\bm B$沿$x$正向，
动生电动势由下式积分：
\[
\mathcal E_{M\to N}
=\int_M^N(\bm v\times\bm B)\cdot d\bm l
=v\int_{a-b}^{a+b}B(x)\,dx
=\boxed{\frac{\mu_0Iv}{2\pi}
\ln\frac{a+b}{a-b}}.
\]
导体上自由电荷达到准静态平衡时，
$\bm E_{\rm static}=-\bm v\times\bm B$，
右端$N$具有较高静电势，因此
\[
\boxed{U_M-U_N
=-\frac{\mu_0Iv}{2\pi}\ln\frac{a+b}{a-b}}.
\]
\crosscheck 结果与半圆弧具体形状无关，
只与杆两端的$x$坐标有关，因为$vB(x)$仅为$x$的函数；
若反向电流或反向速度，则端点电势差换号。

\sourceproblem{9.3.3}{239}{239}
\problemtext 两根平行无限长直导线在同一平面内，
电流方向相反、大小均为$I(t)$，且$\dot I>0$。
在两导线右侧放置矩形线圈，线圈高$\ell$、宽$a$；
从靠右导线到线圈近边的距离为$d$，
从靠左导线到线圈近边距离为$b$（$b>d$）。
求矩形中磁通量及感应电动势。
\workings
采用磁准静态近似，设右边导线的磁场\emph{指向纸内}为正通量（顺时针回路为正），
两根相反电流产生的场分别为
$\mu_0I/(2\pi x)$和反向的场。
对线圈宽度逐条积分，
\[
\Phi_B(t)=\frac{\mu_0I(t)\ell}{2\pi}
\left[
\int_d^{d+a}\frac{dx}{x}
-\int_b^{b+a}\frac{dx}{x}
\right].
\]
因此
\[
\boxed{\Phi_B=
\frac{\mu_0I\ell}{2\pi}
\ln\frac{(d+a)b}{d(b+a)}},
\]
\[
\boxed{\mathcal E_{\rm cw}=
-\frac{\mu_0\ell\dot I}{2\pi}
\ln\frac{(d+a)b}{d(b+a)}}.
\]
这里“cw”表示按纸面向内为正法向、顺时针为正绕行的约定。
\crosscheck 对$b>d$，$\ln[(d+a)b/(d(b+a))]>0$，
说明近导线的贡献较强；
$\dot I>0$时感应电流反抗磁通量增大。

\sourceproblem{9.3.4}{239}{239}
\problemtext 长直导线电流$I=5$ A向上，
右侧同平面矩形线圈高$b=0.06$ m、宽$a=0.04$ m，
近边距导线$d=0.05$ m。
线圈以$v=0.03\,\mathrm{m/s}$向右远离长直导线。
求该瞬间感应电动势的大小和方向。
\workings
导线右侧的磁场指向纸内，
其通过矩形回路的磁通量大小为
\[
\Phi(d)=\frac{\mu_0Ib}{2\pi}
\int_d^{d+a}\frac{dx}{x}
=\frac{\mu_0Ib}{2\pi}\ln\frac{d+a}{d}.
\]
由于$d'(t)=v>0$，
\[
\left|\frac{d\Phi}{dt}\right|
=\frac{\mu_0Ibv}{2\pi}
\left(\frac1d-\frac1{d+a}\right).
\]
代入题设数据，
\[
\boxed{|\mathcal E|
=2.0\times10^{-7}\cdot5\cdot0.06\cdot0.03
\left(20-\frac1{0.09}\right)
=1.60\times10^{-8}\,\mathrm V}.
\]
线圈离导线越来越远，指向纸内的磁通量减小，
为阻止这种变化，感应电流产生\emph{纸内}磁场，
故\textbf{从纸外观察电流为顺时针}。
\crosscheck 面积的宽$a$与高度$b$不能互换：
靠近导线的两个长边$B$不同，
应沿径向宽度积分。

\sourceproblem{9.3.5}{239}{239}
\problemtext 导轨上导体杆$ab$长$\ell$，
以恒定速度$v$向右运动；
磁场$B(t)=kt$（$k>0$），磁感应方向与
导轨平面法线夹角为$60^\circ$。
$t=0$时杆位于固定端$cd$，
求任意时刻闭合回路的感应电动势大小和绕向。
\workings
平面有效法向磁场
$B_\perp(t)=kt\cos60^\circ=kt/2$。
杆从固定端移到$x=vt$，
回路面积$S(t)=\ell vt$，穿过的磁通量
\[
\Phi(t)=B_\perp(t)S(t)
=\frac12k\ell v t^2.
\]
故
\[
\boxed{|\mathcal E(t)|
=\left|\frac{d\Phi}{dt}\right|
=k\ell vt}.
\]
这一总感应电动势可分成磁场自身变化贡献
$k\ell vt/2$与杆运动扩大面积贡献$k\ell vt/2$。
\crosscheck 两项大小\emph{相等}，
漏掉动生或漏掉感生电动势都将得到一半结果。
若从磁场法线正侧观察，正法向通量增长，
由楞次定律感应电流沿\emph{顺时针}方向。

\sourceproblem{9.3.6}{239--240}{239--240}
\problemtext 一边长$a$的方形导线框以恒速$v$向右通过宽度$3a$的均匀磁场区，
磁场方向指向纸内。
$t=0$时方框的右侧前缘刚进入磁场。
以逆时针电流为正，写出整个进出磁场过程中感应电动势的分段表达式。
\workings
在进入阶段$0<t<a/v$，框内与磁场重叠的面积
$S(t)=avt$，指向纸内的磁通量增长，
感应电流沿逆时针，故
$\mathcal E=+Bav$。
当$a/v<t<3a/v$，整个线框位于均匀磁场内，
$S=a^2$恒定，电动势为零。
离开阶段$3a/v<t<4a/v$，重叠面积
$S(t)=a(4a-vt)$，指向纸内的通量减小，
感应电流沿顺时针，故电动势为负。
因此
\[
\boxed{\mathcal E(t)=
\begin{cases}
0,&t<0,\\
+Bav,&0<t<a/v,\\
0,&a/v<t<3a/v,\\
-Bav,&3a/v<t<4a/v,\\
0,&t>4a/v.
\end{cases}}
\]
\crosscheck 进场和离场脉冲持续时间均为$a/v$，
中间零电动势平台宽$2a/v$，
两脉冲积分的有向总和为零，符合初末磁通量同为零。
% === END textbook-ch09-solutions-001-006.tex ===

% === BEGIN textbook-ch09-solutions-007-012.tex blob 5cc6a8835a67d4c30652b2688d5803e663020432 ===
% 《大学物理简明教程》第5版第9章解答题9.3.7—9.3.12
% 源书印刷/PDF第239—240页；核查杆运动几何和电流方向后独立推导。
% 9.3.9对源图几何给出计算值并保留电位正向的约定。

\subsection{解答题：转动导线、涡旋电场与互感}

\sourceproblem{9.3.7}{239--240}{239--240}
\problemtext 长$\ell$的直金属杆$ab$绕过$O$点的垂直轴匀速旋转，
角速度$\omega$；$aO=\ell/3$，$Ob=2\ell/3$，
匀强磁场$\bm B$与旋转轴同向。
求杆两端感应电势差的大小，比较哪端电势高。
\workings
取从$O$向$b$为正的杆坐标$x$，
$x_a=-\ell/3$、$x_b=2\ell/3$。
假设$\bm\omega$与$\bm B$同向，
各点局部动生电动势元为
$d\mathcal E=B\omega x\,dx$。
从$a$到$b$积分：
\[
\mathcal E_{a\to b}
=B\omega\int_{-\ell/3}^{2\ell/3}x\,dx
=\frac{B\omega}{2}
\left[\frac{4\ell^2}{9}-\frac{\ell^2}{9}\right]
=\boxed{\frac{B\omega\ell^2}{6}}.
\]
故
\[
\boxed{U_b-U_a=\frac{B\omega\ell^2}{6}}>0
\]
在该同向轴与电势约定下，\textbf{$b$端电势高}。
\crosscheck 若从图上观察角速度反转或把$\bm B$反向，
电势差有向符号必须反号；
大小为$B|\omega|\ell^2/6$，与杆所处的角位置无关。

\sourceproblem{9.3.8}{240}{240}
\problemtext 非洲某地地磁感应强度$B=4.0\times10^{-5}$ T，
磁场与\emph{水平面}的夹角为$60^\circ$。
将一根长$\ell=1.0$ m、沿东西方向水平放置的金属杆自由下落，
忽略阻力，求$t=3$ s时杆上动生电动势。
\workings
杆沿东西向，速度$\bm v$竖直向下，$v=gt$。
只有磁场的\emph{水平且南北向}分量使
$(\bm v\times\bm B)\cdot d\bm l$沿东西向非零。
已知地磁倾角$60^\circ$，水平场幅为
\[
\boxed{B_{\mathrm H}=B\cos60^\circ=\frac B2}.
\]
若水平场的方位角相对东西方向为$\beta$，
其南北向分量为$B_{\mathrm H}\sin\beta$，
故严谨的一般表达式为
\[
\boxed{|\mathcal E|=B\ell gt\,\cos60^\circ|\sin\beta|}.
\]
若采用通常教材默认的水平地磁场沿南北向
（$\beta=90^\circ$），则
\[
\boxed{|\mathcal E|=4.0\times10^{-5}(1.0)(9.8)(3)
\cdot\frac12
=5.88\times10^{-4}\,\mathrm V}.
\]
\crosscheck \textbf{地磁方位需二校：}源题只给与水平面的倾角，
未明说水平投影方位。在南北向水平分量的常规假设下得
$0.588\,\mathrm{mV}$；其他方位需乘$|\sin\beta|$。
此前将倾角的正弦误作水平分量会导致数值偏大，
本稿按与\emph{水平面}成角的定义使用\emph{余弦}修正。

\sourceproblem{9.3.9}{240}{240}
\problemtext 半径$R$的圆柱形区域内部有均匀、指向纸内的磁场，
$B(t)$随时间增大，圆外磁场为零。
一根长$2R$的水平金属杆$ac$横穿区域下部，
其前半段$ab$（长$R$）在圆内，
后半段$bc$（长$R$）在圆外；
原图半径$Ob$相对竖直方向倾斜$30^\circ$。
求两端感生电动势的大小和沿杆方向。
\workings
令圆心在$(0,0)$，磁场指向纸内，
取杆从$a$到$c$的水平方向为$x$正向。
源图中$Ob$与竖直方向夹角$30^\circ$，
所以水平弦与圆心的距离
$h=\sqrt3R/2$，并取
\[
a=(-R/2,-h),\qquad b=(R/2,-h),\qquad c=(3R/2,-h).
\]
记$\dot B=dB/dt>0$，表示纸内磁场\emph{大小}的增长率。
由轴对称法拉第环流，圆内的感生电场为
\[
\bm E_{\mathrm{in}}
=\frac{\dot B}{2}(-y\,\bm e_x+x\,\bm e_y).
\]
在$y=-h$处，$x$分量恒为$\dot Bh/2$，故
\[
\mathcal E_{a\to b}
=\frac{\dot Bh}{2}R
=\frac{\sqrt3}{4}\dot BR^2.
\]
圆外的感生电场切向分量
$E_\theta=\dot BR^2/(2r)$（沿逆时针），
沿水平线$x$方向的分量为
$E_x=\dot BR^2h/[2(x^2+h^2)]$。
因此
\[
\mathcal E_{b\to c}
=\frac{\dot BR^2}{2}
\left[
\arctan\frac{x}{h}
\right]_{R/2}^{3R/2}
=\frac{\pi}{12}\dot BR^2.
\]
合计
\[
\boxed{\mathcal E_{a\to c}
=\left(\frac{\sqrt3}{4}+\frac{\pi}{12}\right)R^2\dot B}.
\]
感生电场沿杆从\textbf{$a$指向$c$}；
以静电势差记忆时，要说明涡旋电场本身不可写成全局静电势梯度。
\crosscheck 本题必须分段计算\emph{圆内}与\emph{圆外}；
圆外尽管$B=0$，仍有非零涡旋电场。
数值系数基于图中$30^\circ$角确实相对\emph{竖直半径}的解释，
源图方位应二校。

\sourceproblem{9.3.10}{240}{240}
\problemtext 一根无限长直导线载交变电流
$I(t)=I_0\sin\omega t$。
在其右侧共面放一矩形线圈，线圈高$a$（与直导线平行）、
宽$b$，近边与直导线的距离为$d$。
求线圈中的感应电动势。
\workings
在磁准静态近似（忽略传播延迟和辐射，回路尺度远小于电磁波长）下，距直导线$x$处$B(x,t)=\mu_0I(t)/(2\pi x)$，
以纸面向里为正磁通方向，有
\[
\Phi_B(t)=a\int_d^{d+b}B(x,t)\,dx
=\frac{\mu_0aI_0}{2\pi}
\ln\frac{d+b}{d}\sin\omega t.
\]
因此
\[
\boxed{\mathcal E(t)=
-\frac{\mu_0aI_0\omega}{2\pi}
\ln\frac{d+b}{d}\cos\omega t}.
\]
大小峰值
$\boxed{\mathcal E_{\max}
=\mu_0aI_0\omega\ln[(d+b)/d]/(2\pi)}$。
\crosscheck 频率升高会线性提高同一$I_0$条件下的电动势峰值；
固定导线的磁场具有$1/x$的径向分布，不能取为均匀场。

\sourceproblem{9.3.11}{240}{240}
\problemtext 图9.3.11中无限长导线穿过一矩形线圈所在的平面，
长导线电流$I(t)=I_0\sin\omega t$；
矩形高度$a$，位于长导线左侧的宽度$a/2$、
右侧宽度$3a/4$。求互感系数及互感电动势。
\workings
长导线两侧磁场指向纸面的方向相反。
用线圈法向统一定向，两个区域磁通量的有向贡献相减。
在导线半径有限、截面电流轴对称且与线圈电绝缘的模型下，两侧内部磁通对称抵消；对外部贡献，
可把两侧径向积分的公共小半径截止量$\epsilon$
逐项消去：
\[
\Phi_B=\frac{\mu_0Ia}{2\pi}
\left[
\ln\frac{3a/4}{\epsilon}
-\ln\frac{a/2}{\epsilon}
\right]
=\frac{\mu_0Ia}{2\pi}\ln\frac32.
\]
故互感系数大小
\[
\boxed{M=\frac{\mu_0a}{2\pi}\ln\frac32},
\]
互感电动势为
\[
\boxed{\mathcal E(t)=-M\frac{dI}{dt}
=-\frac{\mu_0aI_0\omega}{2\pi}
\ln\frac32\cos\omega t}.
\]
\crosscheck 当左右宽度相等时两侧磁通完全抵消，互感为零；
此处正是两侧宽度不同才出现非零净磁链。
实际长导线有有限半径，以上推导中两边同一截止尺度相消。若理想化为零半径线电流，这一积分是对称柯西主值而非两个独立收敛的积分；若导线与线圈电接触或回流路径不可忽略，应重新建模。

\sourceproblem{9.3.12}{240}{240}
\problemtext 一对相互平行的长直导线，其横截面圆半径同为$a$，
两导线中心相距$d$，组成一往返电路。
忽略导线\emph{内部}磁通量，证明每段长度$\ell$的自感为
$L=(\mu_0\ell/\pi)\ln[(d-a)/a]$。
\workings
两导线载电流大小$I$、方向相反，
对两导线间跨径向的狭长曲面，
两导线磁场沿法向\emph{同向}。
从靠近第一根导线表面$r=a$积分到
靠近第二根表面$r=d-a$：
\[
\Phi_B=\frac{\mu_0 I\ell}{2\pi}
\int_a^{d-a}\left(\frac1r+\frac1{d-r}\right)dr
=\frac{\mu_0I\ell}{\pi}\ln\frac{d-a}{a}.
\]
因此
\[
\boxed{L=\frac{\Phi_B}{I}
=\frac{\mu_0\ell}{\pi}\ln\frac{d-a}{a}}.
\]
\crosscheck 需要$d>2a$使两圆柱互不相交；
该式是外部磁链的几何近似，
不包含两导线内部的电流分布所带来的内部自感。
% === END textbook-ch09-solutions-007-012.tex ===

% === BEGIN textbook-ch09-solutions-013-017.tex blob d9444d8cde57a1dac5a6cd91fd75a96c28c00ea2 ===
% 《大学物理简明教程》第5版第9章解答题9.3.13—9.3.17
% 原书/PDF第240—241页；依据扫描图示，独立求解磁场自能和位移电流。

\subsection{解答题：螺绕环、自感磁能与位移电流}

\sourceproblem{9.3.13}{240--241}{240--241}
\problemtext 一矩形截面的空心螺绕环共$N$匝，
横截面高度$h$、内径向半径$a$、外径向半径$b>a$，
空气芯，线圈电流$I$。求该线圈的自感和储存的磁能。
\workings
理想均匀密绕的环形线圈内部，
由安培定理
\[
\bm B(r)=\frac{\mu_0NI}{2\pi r}\bm e_\theta,
\qquad a<r<b.
\]
一匝所链接的截面磁通为
\[
\Phi_{\mathrm{one}}=\int_a^b B(r)h\,dr
=\frac{\mu_0NIh}{2\pi}\ln\frac ba.
\]
总磁链$\Psi=N\Phi_{\mathrm{one}}$，故
\[
\boxed{L=\frac{\Psi}{I}
=\frac{\mu_0N^2h}{2\pi}\ln\frac ba},
\]
\[
\boxed{W_m=\frac12 LI^2
=\frac{\mu_0N^2hI^2}{4\pi}\ln\frac ba}.
\]
\crosscheck $L$与匝数平方成正比、与电流无关；
当$b\to a^+$时截面趋零，自感趋于零，
符合理想环形线圈模型。

\sourceproblem{9.3.14}{241}{241}
\problemtext 无限长实心圆柱导体半径$R$，总电流$I$
均匀分布于横截面，求\emph{导体内部}每单位长度存储的磁场能。
\workings
半径$r<R$的同轴圆形安培环路包围
$I_{\mathrm{enc}}=I r^2/R^2$，
所以
$B(r)=\mu_0Ir/(2\pi R^2)$。
真空磁能密度为$w_m=B^2/(2\mu_0)$。
取单位轴向长度，体积元$dV=2\pi r\,dr$，
\[
\frac{W_{\mathrm{inside}}}{\ell}
=\int_0^R\frac{B(r)^2}{2\mu_0}\,2\pi r\,dr
=\frac{\mu_0I^2}{4\pi R^4}\int_0^Rr^3\,dr.
\]
故
\[
\boxed{\frac{W_{\mathrm{inside}}}{\ell}
=\frac{\mu_0I^2}{16\pi}}.
\]
\crosscheck 在电流固定且均匀分布下，结果不依赖半径$R$；
若把\emph{导体外部}也算入无限长孤立导线的全部空间，
还会有发散的对数外场能，所以本题特别限于导体内部。

\sourceproblem{9.3.15}{241}{241}
\problemtext 长圆柱形电容器内外圆柱电极半径为
$R_1<R_2$，中间填充介电常数$\varepsilon$
（SI绝对介电常数）的均匀线性介质。
施加的两板电压随时间满足$dU/dt=k$（常数）。
求距轴线$r$、$R_1<r<R_2$处的位移电流密度。
\workings
圆柱轴对称电准静态近似、无介质漏电且$U=\varphi_{\rm inner}-\varphi_{\rm outer}$时，
径向电场大小为
\[
E_r(r,t)=\frac{U(t)}{r\ln(R_2/R_1)}.
\]
介电位移场$\bm D=\varepsilon\bm E$，
其时间导数给出位移电流密度：
\[
\boxed{\bm J_D(r)=\frac{\partial\bm D}{\partial t}
=\frac{\varepsilon k}{r\ln(R_2/R_1)}\bm e_r},
\quad R_1<r<R_2,
\]
其中径向正向定义为内极板指向外极板。
\crosscheck 对同轴圆柱面面积$2\pi r\ell$积分，
$I_D=2\pi\ell\varepsilon k/\ln(R_2/R_1)$，
不随$r$改变，符合同一电容器电流连续性。

\sourceproblem{9.3.16}{241}{241}
\problemtext 证明平行板电容器在板距、面积和介质不随时间改变时，
位移电流满足$I_D=C\,dU/dt$；
讨论这一关系能否推广到非平行板电容器。
\workings
平行板电容器忽略边缘效应，
$C=\varepsilon S/d$、
$E=U/d$、$D=\varepsilon E$，因此
\[
I_D=\frac d{dt}\int_S\bm D\cdot d\bm S
=\varepsilon S\frac{d}{dt}\frac Ud
=\boxed{C\frac{dU}{dt}}.
\]
对于\emph{几何固定、线性且介电参数不随时间改变}
的其他二端电容器，也有端电荷$Q=CU$，因而
\[
\boxed{I_D=\frac{dQ}{dt}=C\,\frac{dU}{dt}}
\]
适用于端口电流连续意义下的相应截面。
\crosscheck 不能无限制推广：
对于固定的绝缘介质截面，上述端口关系成立；若有传导漏电，端电流还须计入传导电流。若几何或介电性质随时间改变使$C=C(t)$，
应写$I=\frac{d(CU)}{dt}=C\dot U+U\dot C$；
非线性介质还须按真实$Q(U)$微分。

\sourceproblem{9.3.17}{241}{241}
\problemtext 真空中半径$R=0.10$ m的圆形平行板电容器
均匀充电，板间电场增长率为
$dE/dt=1.0\times10^{13}\,\mathrm{V/(m\,s)}$。
求板间位移电流大小、板中心轴线距离$r<R$处的磁感应强度
$B_r$，以及$r=R$处的磁感应强度$B_R$。
\workings
忽略边缘效应，位移电流为
\[
I_D=\varepsilon_0\,\frac{dE}{dt}\,\pi R^2.
\]
取$\varepsilon_0=8.8541878128\times10^{-12}\,\mathrm{F/m}$，
得到
\[
\boxed{I_D=
\varepsilon_0\pi(0.10)^2(10^{13})
\approx2.78\,\mathrm A}.
\]
在两板间取半径$r<R$的圆形安培环路，
包围位移电流
$I_D(r)=\varepsilon_0\pi r^2\,dE/dt$，
由安培--麦克斯韦方程
$B(r)(2\pi r)=\mu_0I_D(r)$，故
\[
\boxed{B_r=\frac{\mu_0\varepsilon_0r}{2}
\frac{dE}{dt}},\qquad 0\le r<R.
\]
在板边缘
\[
\boxed{B_R=\frac{\mu_0\varepsilon_0R}{2}
\frac{dE}{dt}
\approx5.56\times10^{-6}\,\mathrm T}.
\]
\crosscheck 对$r>R$，
若忽略边缘漏场，环路所包的位移电流不再随$r$增加，
有$B(r)\propto1/r$；
内外解在$r=R$处连续。
% === END textbook-ch09-solutions-013-017.tex ===
